\titolosezione{Il fuoco di bivacco}
\begingroup\small
Un buon fuoco ha un andamento a onda: sale, tocca il punto più alto e poi scende fino al silenzio.
Scegliere i canti seguendo l'intensità (da \fiamme{1} a \fiamme{5}) aiuta a non bruciare tutto subito.

\medskip
\begin{tabular}{@{}l@{\hspace{4mm}}p{0.72\textwidth}@{}}
\textbf{Apertura} & \fiamme{3}\ Si accende il fuoco, ci si riunisce. Canti noti, facili da attaccare, che fanno sentire tutti dentro.\\[1.5mm]
\textbf{Salita} & \fiamme{4}\ Bans, canoni e canti con gesti: si scioglie il gruppo e si prende ritmo.\\[1.5mm]
\textbf{Picco} & \fiamme{5}\ Il momento più acceso, breve. Canti gridati, balli, voci che si sovrappongono.\\[1.5mm]
\textbf{Discesa} & \fiamme{3}\ Si abbassa piano la voce e il ritmo; si torna a cantare insieme, non uno contro l'altro.\\[1.5mm]
\textbf{Raccoglimento} & \fiamme{2}\ Canti lenti, di strada e di ricordo. Il fuoco è basso, si canta quasi sottovoce.\\[1.5mm]
\textbf{Chiusura} & \fiamme{1}\ Un canto semplice o una preghiera, poi silenzio. Non si applaude: si lascia che resti.
\end{tabular}

\medskip
\textit{Regola d'oro:} chi guida legge il fuoco, non la scaletta. Se il gruppo è stanco, si scende prima.
\endgroup
